\documentclass[a4paper,10pt]{article}
\usepackage[utf8]{inputenc}
\usepackage{graphicx}
\usepackage{subfigure}
\usepackage{fancyhdr}
\usepackage{listings}
\usepackage{ gensymb }


%opening

\pagestyle{fancy} %eigener Seitenstil
\fancyhf{} %alle Kopf- und Fußzeilenfelder bereinigen
\fancyhead[L]{SEP Abgabegespräch SS2015/16} %Kopfzeile links
\fancyhead[C]{} %zentrierte Kopfzeile
\fancyhead[R]{\includegraphics[width=1cm]{../../template/tugraz_bw}} %Kopfzeile rechts
\renewcommand{\headrulewidth}{0.4pt} %obere Trennlinie
\fancyfoot[C]{} %Seitennummer
\renewcommand{\footrulewidth}{0.4pt} %untere Trennlinie


\lstset{language=C++, basicstyle=\ttfamily\tiny}  

\begin{document}

\section*{Vokabeltrainer}

Schreibe ein einfaches Programm, mit welchem Vokabeln gelernt werden können. Dabei liest das Programm die Vokabeln aus einer Wörterbuchdatei ein und der Benutzer soll die korrekte Übersetzung des Wortes eingeben. Die Wörterbuchdatei hat folgendes Format: \texttt{<Englisches Wort> <Deutsches Wort>}.\newline
Beispiel für eine Liste \texttt{vocabulary.txt}:
\begin{lstlisting}[frame=single]
hello hallo
house haus
bell klingel
pen stift
cat katze
\end{lstlisting}

Folgende Befehle stehen zur Verfügung
\begin{description}
 \item[quit]: Beendet das Programm.
 \item[next $<$language$>$]: Gibt das nächste Wort in der gewählten Sprache aus.
 \item[answer $<$word$>$]: Lösung des Benutzers.
 \item[help]: Lösung wird ausgegeben, wird als falsch bewertet.
 \item[show progress]: Zeigt die Anzahl der verbleibenden Wörter und die Anzahl der richtigen und falschen Wörter.
\end{description}
Zu beachten ist, dass die Reihenfolge der Wörter zufällig sein soll, ein erneuter Programmaufruf mit der gleichen Wörterbuchdatei soll eine neue Reihenfolge haben.

\subsection*{Beispiel}
\footnotesize
\begin{verbatim}
./trainer vocabulary.txt
> next german
hallo
> answer hello
Correct answer!
> next english
pen
> answer hund
Wrong answer!
> next german
klingel
> help
answer is: bell
> show progress
4 of 5 words are answered. 1 correct, 2 wrong
> quit
\end{verbatim}

\subsection*{Bewertung}
\begin{itemize}
 \item Dokumentation und Programmierstil: 1 Punkt
 \item Einhalten der OOP Konzepte: 1 Punkt
 \item Struktur \& Korrektheit: 1 Punkt
\end{itemize}
\newpage


\end{document}
